\section{Multi-Head Self-Attention}\label{sec:multi_self_attention}
\textit{Describe the working of multi-head self-attention. Write PyTorch code to implement
scaled dot-product attention and extend it to multi-head attention. Explain how
multiple heads improve representation learning.}

Multi-head self-attention is the central component of the
Transformer architecture and is responsible for modeling
dependencies between tokens within a sequence.
Unlike recurrent models, self-attention allows every token
to directly attend to all other tokens in the sequence,
enabling efficient capture of both short-range and long-range
relationships.
This mechanism forms the basis of the Transformer model
introduced in~\cite{vaswaniAttentionAllYou2017}.

Self-attention operates by projecting an input sequence into
three distinct representations: queries, keys, and values.
Given an input matrix
$X \in \mathbb{R}^{T \times d_{\text{model}}}$, these projections
are computed using learned weight matrices
\begin{equation}
Q = XW_Q, \quad K = XW_K, \quad V = XW_V,
\end{equation}
where $W_Q$, $W_K$, and $W_V$ are trainable parameters.
The attention mechanism then computes pairwise similarity
scores between queries and keys to determine how much
each token should attend to every other token.

The core operation is known as \emph{scaled dot-product
attention}, which is defined as
\begin{equation}
\text{Attention}(Q, K, V) =
\text{softmax}\left(\frac{QK^\top}{\sqrt{d_k}}\right)V,
\end{equation}
where $d_k$ is the dimensionality of the key vectors.
The scaling factor $\sqrt{d_k}$ prevents the dot products
from growing too large, which stabilizes training by
keeping gradients well-behaved.
The softmax operation ensures that the attention weights
form a valid probability distribution.

While a single attention operation is powerful, using
multiple attention heads significantly improves the
expressiveness of the model.
Multi-head attention splits the input representations
into $h$ smaller subspaces and performs attention
independently in each subspace.
Each head attends to the sequence using different learned
projections, allowing the model to capture diverse
relationships such as syntactic structure, semantic
similarity, or positional alignment.

Formally, for $h$ attention heads, the multi-head attention
operation is defined as
\begin{equation}
\text{MultiHead}(Q, K, V) =
\text{Concat}(\text{head}_1, \dots, \text{head}_h)W_O,
\end{equation}
where each head is computed as
\begin{equation}
\text{head}_i = \text{Attention}(Q_i, K_i, V_i),
\end{equation}
and $W_O$ is a learned output projection matrix.
By concatenating the outputs of multiple heads, the model
combines information from different representation
subspaces into a single unified embedding.

A PyTorch implementation of scaled dot-product attention
and its extension to multi-head self-attention is shown in
\Cref{lst:multihead_self_attention}.
The code illustrates how queries, keys, and values are
reshaped to support multiple heads, how attention is
computed in parallel, and how the outputs are concatenated
and projected back to the original embedding dimension.
\begin{listing}
    \begin{minted}{Python}
class MultiHeadSelfAttention(nn.Module):
    def __init__(self, d_model, num_heads):
        super().__init__()
        assert d_model % num_heads == 0

        self.num_heads = num_heads
        self.d_k = d_model // num_heads

        self.W_q = nn.Linear(d_model, d_model)
        self.W_k = nn.Linear(d_model, d_model)
        self.W_v = nn.Linear(d_model, d_model)
        self.W_o = nn.Linear(d_model, d_model)

    def scaled_dot_product_attention(self, Q, K, V):
        # Q, K, V: (batch, heads, seq_len, d_k)
        scores = torch.matmul(Q, K.transpose(-2, -1)) / (self.d_k ** 0.5)
        weights = torch.softmax(scores, dim=-1)
        return torch.matmul(weights, V)

    def forward(self, x):
        # x: (batch, seq_len, d_model)
        B, T, _ = x.size()

        Q = self.W_q(x).view(B, T, self.num_heads, self.d_k).transpose(1, 2)
        K = self.W_k(x).view(B, T, self.num_heads, self.d_k).transpose(1, 2)
        V = self.W_v(x).view(B, T, self.num_heads, self.d_k).transpose(1, 2)

        attn = self.scaled_dot_product_attention(Q, K, V)

        attn = attn.transpose(1, 2).contiguous().view(B, T, -1)
        return self.W_o(attn)
    \end{minted}
    \caption{Minimal Implementation of Multi-Head Self-Attention}
    \label{lst:multihead_self_attention}
\end{listing}

